%\section{n-fold loop spaces}
%Operads were defined by many people; Artamonov ~'69, Boardman-Vogt '68, May '72, also Lambek '69 in the form of multicategories.

%For homotopy theorists, maybe the most influential is May, who wanted to study the homotopy types of n-fold loop spaces. He proved that they could be identified by their algebraic properties. We notice that in a loop space we can define concatenation operations $\gamma \bullet \delta$ uniquely up to homotopy, and while this is not associative it is associative up to homotopy, and furthermore these homotopies are not associative on the nose but are up to higher homotopies, etc. In other words $\Omega X$ is a homotopy-coherent version of a monoid, called an $A_\infty=E_1$ space. If you have a double loop space $\Omega^2 X$ then you have two choices of concatenation; external and internal. By an Eckmann-Hilton style argument these correspond to a single operation commutative up to homotopy. However, these homotopies, and their higher homotopies, don't commute - this is called an $E_n$ space. May's recognition principle then states: 
%\begin{theorem}[May]
%    A connected space $Y$ is weakly homotopy equivalent to $\Omega^nX$ if and only if $Y$ is an $E_n$-space.
%\end{theorem}

%To make this statement clear you need to be able to control these new variants of classical algebraic structures. This leads us to operads.

\begin{comment}
    
\section{Operads and multicategories}
Operads are a combinatorial gadget for controlling an algebraic theory in a monoidal category. They are a special case of a more general construction, called a coloured operad or multicategory. A multicategory is a generalization of a category where morphisms are allowed to have multiple inputs.

The notion is very similar to that of a Lavwere theory, but operads can be evaluated in any symmetric monoidal category.

\begin{definition} A (symmetric) multicategory or coloured operad $\mathcal{O}$ is
    \begin{itemize}
        \item A collection $ob(C)$ of objects
        \item For each natural number $n\geq 0$, $(X_i)_{1\leq i \leq n}$ and $Y$ a hom-set $Hom(X_1,\dots,X_i;Y)$
        \item An identity $id\in Hom(X;X)$
        \item A multicomposition operation $$\prod_j Hom((X_{ij})_i;Y)\times Hom(Y_1,\dots, Y_j;Z)\rightarrow Hom((X_{ij})_{ij};Z)$$
        \item For each $\sigma\in S_n$ a map $\sigma^*:Hom(\{X_i\}_{1\leq i \leq n};Y)\rightarrow Hom(\{X_{\sigma(i)}\}_{1\leq i \leq n};Y)$ forming a representation of $S_n.$
    \end{itemize}
    Such that the composition is unital, associative and $S_n$-equivariant.
\end{definition}

    The $S_n$ action tells us what happens when we permute inputs.
 $S_n-$equivariance means $$\gamma\sigma\circ(f^{k_1}_1,\dots,f^{k_n}_n)=\big(\gamma\circ(f_{\sigma^{-1}(i)}^{k_{\sigma^{-1}(i)}})_i\big)\sigma(k_1,\dots,k_n)$$ (draw a picture)
    
    It is evident how to generalize this to enriched variants.



    \begin{example}
        A symmetric monoidal category.
    \end{example}
    
    \begin{definition}
    A symmetric operad in $V$ is exactly a single object $V-$enriched symmetric multicategory, suitably defined. You may see me draw objects as multimorphisms.
    \end{definition}



%Operads are a special case of (pointed, symmetric) multicategories which I now define.

\begin{definition} A symmetric operad in $(Set,\times,1)$ is a one-object symmetric multicategory.
    \item Composition is equivariant with respect to the symmetric action. In other words, $$\gamma(f\sigma;g_1,\dots,g_k)=\gamma(f;g_{\sigma^{-1}(1)},\dots,g_{\sigma^{-1}(k)})\sigma(j_1,\dots,j_k)$$ where $g_i\in \mathcal{O}(j_i)$, $f\in \mathcal{O}(k)$ and $\sigma\in\Sigma_k$, as well as a similar condition when starting with actions $\tau_1,\dots,\tau_k$. In essence, we know how to "permute the elements" and this acts sensibly with the multicomposition.
\end{definition}


\begin{example}
    For each $X\in C$ monoidal, the \defn{endomorphism operad} $End(X)$ has as $n^{th}$ object $C(X^{\tensor n},X)$, with the unit $e$ and obvious composition morphisms.
\end{example}
\begin{remark}
Composition in the operad is associative, but the \textit{operations themselves} need not be associative. For example, in the endomorphism operad, there is nothing requiring that an arbitrary $f:X\tensor X\rightarrow X$ satisfies $f(f(x,y),z)=f(x,f(y,z))$. 
\end{remark}

\begin{definition}
    A morphism of operads $\mathcal{O}$ and $\mathcal{U}$ is a collection of $\Sigma_n$-invariant morphisms $\mathcal{O}(n)\rightarrow \mathcal{U}(n)$ for each $n$ that respects composition of each operad, i.e. so the following diagram commutes for each choice of number of inputs.
    % https://tikzcd.yichuanshen.de/#N4Igdg9gJgpgziAXAbVABwnAlgFyxMJZABgBpiBdUkANwEMAbAVxiRAB12BbOnACwDGjYAHkAvgAowASk4Q8XeAAJOPfkIajJAawD6ARlnt5WRXE5R554wuWreg4eIl6ZIMaXSZc+QijL6VLSMLGz26sIAqpIycrZwKtwOGsDRLgZGJmYWVnGmdkkRmmmu0u6eIBjYeARE+uRB9MysiByFjprO2mUeXtW+daSB1E2hreEdqTo9QTBQAObwRKAAZgBOEFxIZCA4EEj6vSDrm9vUe0gATEcnW4j1u-uIAMw3G3eX50+vFGJAA
\[\begin{tikzcd}
\mathcal{O}(n)\otimes \mathcal{O}(k_1)\otimes\dots\otimes \mathcal{O}(k_n) \arrow[d] \arrow[r] & \mathcal{O}(k) \arrow[d] \\
\mathcal{U}(n)\otimes \mathcal{U}(k_1)\otimes\dots\otimes \mathcal{U}(k_n) \arrow[r]           & \mathcal{U}(k)          
\end{tikzcd}\]
\end{definition}

\begin{definition}
    An \defn{algebra} over an operad $\mathcal{O}$ is an object $X$ and a morphism of operads $\theta:\mathcal{O}\rightarrow End(X)$. A \defn{morphism of algebras} from $\theta$ to $\theta'$ is a natural transformation $\theta\Rightarrow \theta'$.
\end{definition}

\begin{example}
    Consider $(Set,\times)$. Let $Ass$ be the operad where $Ass(n)=\Sigma_n$ as a set, with the induced right action of $\Sigma_n$. Then an $Ass-$algebra is a set $X$ along with $n!$ maps $X^n\rightarrow X$. You should think of these as the $n!$ ways to permute $n$ objects. By the condition on the symmetric action, these can be built from each other by permuting inputs. In particular, by the composition law, all operations are built from a single operation $X^2\rightarrow X$, which is unital and associative (since $\Sigma$ is). An $Ass$-algebra is therefore the same as a monoid.

    Similarly, $Com$ is the operad where $Com(n)$ is the singleton set. An algebra over $Com(n)$ is therefore a unique map $X^n\rightarrow X$ for each $n$ that is invariant under permutation of elements, unital, and built from the map $X^2\rightarrow X$. This is the same as a commutative monoid.
\end{example}

\section{$\infty-$categories}
An $\infty-$category is a homotopy-coherent version of a category. I invite you to think of any of the following:
\begin{itemize}
    \item A class of objects, 1-morphisms between objects, 2-morphisms between 1-morphisms, etc.. Such that composition is defined only up to a higher morphism, and all n-morphisms for $n\geq 2$ are invertible.
    \item A category enriched in Top, or Kan.
    \item A simplicial set with inner horn fillers.
\end{itemize}

An $\infty-$category is the natural place to do mathematics in any setting that comes with a notion of homotopy.

All* the standard constructions and theorems of category theory apply in the $\infty-$categorical context, so when working with them you can just think of them as ordinary categories, except every statement has to be formulated in terms of a universal property - you are not allowed to mention the words "element", "equality", "unique", etc.

A standard category can canonically be considered an $\infty-$category and every definition will specialize to one you know. You're more than welcome to only think about 1-categories.


\section{$\infty$-operads}
We really mean $\infty-$multicategories.

If homotopy theory has taught us anything it is that homotopical objects are better behaved when we keep track of the homotopies. We should therefore really be seeking a notion of operad intrinsic to $\infty-$category theory. The standard process is this: take a definition from $1-$category theory, extend it losslessly to its nerve, then extend that definition to all $\infty-$categories.

A good strategy is to rewrite your definition as a lifting property, as we will now see.

Really, we will generalize the notion of symmetric multicategory, or coloured operad.
\section{Symmetric monoidal $\infty-$categories}
To a symmetric monoidal category $(C,\tensor,1)$ we can associate a category $C^\tensor$ where
\begin{itemize}
    \item objects are finite sequences $[C_1,\dots,C_n]$ of objects in $C$,
    \item A morphism from $[C_1,\dots,C_n]$ to $[C_1',\dots,C_m']$ is a subset $S\subset \langle n \rangle^\circ$, a map $\alpha:S\rightarrow \langle m\rangle^\circ\}$ and a collection of maps $$\{f_j:\tensor_{\alpha(i)=j}C_i\rightarrow C_j'\}_j$$ where we define $1:=\tensor_{\emptyset}$.
    \item Composition is the obvious composition of multimorphisms.
\end{itemize}

There is a forgetful functor $p:C^\tensor\rightarrow Fin_*$ which satisfies two properties:
\begin{itemize}
    \item $p$ is a coCartesian fibration (op-fibration): for every $p(f):\langle m\rangle \rightarrow \langle n \rangle$ and $C=[C_1,\dots,C_m]$ there is a cocartesian lift $f:C\rightarrow C'$ over $f$. That is, $f$ has the following universal property for every $g:C\rightarrow D$.
% https://tikzcd.yichuanshen.de/#N4Igdg9gJgpgziAXAbVABwnAlgFyxMJZABgBoBGAXVJADcBDAGwFcYkQBhEAX1PU1z5CKchWp0mrdhwDkPPiAzY8BImWLiGLNohAARef2VC1pAMybJOkAB0bjemADmjGAAIwbuwCdHLtrxGgqoi5pba7HYOzq5uALZeNr4xAQpKwcIkpABM4VK6UX6xANaJyf484jBQTvBEoABm3hBxSGQgOBBIohIRug2GIE0tbTSdSNk0jFhg1nAQ01AgNFr5IE7LIA4ARjCMAAoCKsJbMA04g8OtiGZjXYgALCtW7GgAFA0AlJs7e4fGIRA3iwTgAFhdAkNmtceuNEJNems7DAAB5YOA4OBuADucim9F2ByOJl0wLBFymM2sUAgOBw1Uu0KQTw69wArM8+rYbA0IL5GIwcT8CX9iYCyeDGSMbnckBzEdZ3k5vtxKNwgA
\[\begin{tikzcd}
D                                                      &                                            \\
C \arrow[r, "f"] \arrow[u, "g"]                        & C' \arrow[lu, "\exists w'"', dotted]       \\
\langle k \rangle                                      &                                            \\
\langle n \rangle \arrow[r, "p(f)"'] \arrow[u, "p(g)"] & \langle m \rangle \arrow[lu, "\forall w"']
\end{tikzcd}\]
This cocartesian morphism is the unique morphism whose components are isomorphisms.

Equivalently, every morphism $\langle n \rangle\rightarrow \langle m \rangle$ gives rise to a functor of fibres $C^\tensor_{\langle n \rangle}\rightarrow C^\tensor_{\langle m \rangle}.$ I.e. describes a functor $FinSet_*\rightarrow Cat$. This is straightening-unstraightening.
    \item $C^\tensor_{\langle n \rangle}\simeq C^n$ induced by the equivalence you're thinking about.
\end{itemize}

Here the fiber $C_{\langle n \rangle}^\tensor$ is the full subcategory that lies over $id:\langle n \rangle \rightarrow \langle n\rangle$. Note in particular this implies that every $f:\langle n \rangle \rightarrow \langle m \rangle$ induces a functor of fibers $F_f:C^\tensor_{\langle n\rangle}\rightarrow C^\tensor_{\langle m\rangle}$ given on objects by the codomain of the cocartesian morphism $f'$, and on morphisms given by the unique $F_f(g)$ associated to $g:C_1\rightarrow C_2$ as in the diagram below.

% https://tikzcd.yichuanshen.de/#N4Igdg9gJgpgziAXAbVABwnAlgFyxMJZABgBoBmAXVJADcBDAGwFcYkQAdDx+sAc0YwABGCFcATrwHCQAX1LpMufIRQBGCtTpNW7Lj36ChAWzEdJhtvMXY8BImQBMWhizaJO3KUdETvVhRAMWxUiDWcaV10PfX8TMwtpOUDg5XsUMjUXHXcQAGEAfTVkmzTVZA0syJz2ABEikqClO3KyYmy3dkLHRtSWsNJ26s6Pep7ZLRgoPngiUAAzcQhjJDIQHAgkDW0RkHnGxeWkABYaDaQAVmHovYByA6WVxFP1zcQANmvcvgejxDXzohHF92FgoAVgLFLCIEv5ZL8np9XkgAOwgjzzW73ax7R5IYHIxDkdF7BFbM5vYk7G5giFQ6TxPyWeE0HgAIxgjAACs1Qh5xFg+AALHBkxBXQlo6m5ABiAAo+ABKECs+gc7m89IgAXC0WsrBgXJQehwIVTOSUWRAA
\[\begin{tikzcd}
C_2 \arrow[r, "f''"]                                                  & D_2                                                    \\
C_1 \arrow[r, "f'"] \arrow[u, "g"]                                    & D_1 \arrow[u, "F_f(g)"', dashed]                         \\
\langle n \rangle \arrow[r, "f"]                                      & \langle m \rangle                                      \\
\langle n \rangle  \arrow[r, "f"] \arrow[u, "id_{\langle n \rangle}"] & \langle m \rangle \arrow[u, "id_{\langle m \rangle}"']
\end{tikzcd}\]
Functoriality is ensured by the uniqueness of this lift.

The symmetric monoidal category $(C,\tensor, 1)$ is uniquely determined by a choice of category $C^\tensor$ and functor $p:C^\tensor \rightarrow Fin_*$ satisfying these two conditions. Indeed suppose $p:D\rightarrow Fin_*$ is such a functor. Then we build a symmetric monoidal category $C$ such that
\begin{itemize}
    \item $C=D_{\langle 1 \rangle}$. Note this is the fiber over the identity morphism.
    \item The unit is the codomain of the cocartesian lift over $\langle 0\rangle \rightarrow \langle 1 \rangle $. Note a map $[]\rightarrow [C]$ is exactly a map $1\rightarrow C$. $1$ is initial in $C/1.$
    \item The monoidal product is given by the cocartesian lift over the map $\langle 2\rangle \rightarrow \langle 1\rangle$ $1,2\mapsto 1$ alongwith the isomorphism $C\times C \iso D_{\langle 2 \rangle}$.
    \item The cocartesian morphism over $\langle 1\rangle \rightarrow \langle 2\rangle \quad 1\mapsto 1$ starting at $C$ is necessarily $C\mapsto [C,1]$, and since compositions of cocartesian morphisms are cocartesian, and the cocartesian morphism over $id_{\langle 1\rangle}$ is $id_C$, this implies $C\tensor 1\iso C.$
    \item Considering the two bijections $\langle 2\rangle \rightarrow \langle 2 \rangle$ gives the symmetry of the monoidal product.
    \item Associativity on the monoidal product is given by considering the two maps $\langle 3\rangle \rightarrow \langle 2 \rangle$ given by
    $$(1,2\mapsto 1,3\mapsto 2) \text{  and  } (1\mapsto 1, 2,3\mapsto 2)$$
    and noting that they are equalized by the map $\langle 2 \rangle \rightarrow \langle 1\rangle$ defining the monoidal product.
\end{itemize}

These remarks allow us to define symmetric monoidal $\infty$-categories as follows:
\begin{definition}
    A \defn{symmetric monoidal $\infty$-category} is an $\infty-$category $C^\tensor$ and a coCartesian fibration\footnote{A coCartesian fibration is an inner fibration with all p-coCartesian lifts} $C^\tensor \rightarrow \text{FinSet}_*$ that furthermore defines equivalences $C^\tensor_{\langle n \rangle}\iso (C^\tensor_{\langle 1 \rangle})^n.$
\end{definition}

As in the $1-$categorical world, any $\infty-$category admitting finite products (respectively coproducts) can be promoted to a symmetric monoidal $\infty-$category.
\end{comment}

\begin{comment}

\section{$\infty$-operads}

\begin{definition}
    An \defn{inert-active} factorisation system on a category $C$ is two wide subcategories $(I,A)$ that both contain all isomorphisms, such that any morphism $f\in Mor(C)$ factorizes uniquely up to unique iso by an inert followed by an active morphism.
\end{definition}

\begin{example}
    $FinSet_*$ has an inert-active factorisation system where the inerts are "bin maps" $f:\langle n\rangle\rightarrow \langle m \rangle$ such that $f^{-1}(\langle m \rangle^\circ)\rightarrow \langle m \rangle^\circ$ is an iso, and the active maps are $f$ such that $f^{-1}(*)=\{*\}.$

    Define $p^i:\langle n \rangle \rightarrow \langle 1 \rangle$ as the unique inert with $p^i(i)=1.$
\end{example}

If we relax the definition of a symmetric monoidal $\infty-$category, we can arrive at a definition of $\infty-$operads.

\begin{definition}
    An \defn{$\infty-$operad} is a (categorical fibration\footnote{equivalences lift to equivalences. The axioms actually imply $p$ is a cat. fib. if $p$ is any such functor.}) $p:\mathcal{O}^\tensor \rightarrow FinSet_*$ such that
    \begin{enumerate}
        \item Inert morphisms admit p-coCartesian lifts.
        \item For each $f:\langle m \rangle \rightarrow \langle n \rangle$ and $p-$coCart mors $C'\rightarrow C_i'$ lying over $\rho^i:\langle n \rangle \rightarrow\langle 1 \rangle$ for $1\leq i \leq n$ the induced map 
        $$Map_{O^\tensor}^f(C,C')\rightarrow\prod_i Map_{O^\tensor}^{\rho^i \circ f}(C,C_i')$$
        is a homotopy equivalence.
        \item For every $C_1,\dots,C_i\in \mathcal{O}:=\mathcal{O}^\tensor_{\langle 1 \rangle}$ there is a $C\in \mathcal{O}^\tensor_{\langle n \rangle}$ and p-coCartesian morphisms $C\rightarrow C_i$ covering $p^i:\langle n \rangle \rightarrow \langle 1 \rangle$. Alternatively, $\mathcal{O}^\tensor_{\langle n \rangle}\iso \mathcal{O}^n.$
    \end{enumerate}
\end{definition}

2. says the maps in the category of operators is determined by its projection to a single point. I.e. we don't need a notion $Mul(n,m)$.

\begin{remark}
    We can recover the more pedestrian notion of colored operads as follows:
    \begin{itemize}
        \item The space of objects $\mathcal{O}$ is just the fiber over $\langle 1\rangle$.
        \item (1) in particular implies that inert $\alpha$ induce functors $\mathcal{O}^\tensor_{\langle n \rangle}\rightarrow  \mathcal{O}^\tensor_{\langle m \rangle}$. 
        \item Identities are $p-$cocartesian lifts of the inert map $\langle 0 \rangle \rightarrow \langle 1 \rangle$
        \item Multimorphism spaces $Mul(\{X_i\},Y)$ are the space of morphisms over the active map $\langle n \rangle \rightarrow \langle 1 \rangle$.
        \item Composition is (?)
    \end{itemize}
\end{remark}

%    Here is a pedestrian translation of each of these axioms -- let's focus on the case of monocolored operads. In this case, that corresponds to furthermore equipping $\mathcal{O}^\tensor$ with a homotopy equivalence $\Delta^0\rightarrow \mathcal{O}$

%    First note $p-$coCartesian morphisms 
%    \begin{itemize}
%        \item 1 in particular says that $\rho^i$ admits p-coCartesian lifts.
%        \item p-coCartesian morphisms over $\rho^i$ correspond to multiplication sets $Mul({C_i},C')$.  
%        \item 3 says that objects can be thought of as sequences of objects in $\mathcal{O}$
%        \item 2 then says that morphisms over $\alpha$ are given by $$Hom(\{X_i\},\{Y_j\})\iso \{Mul(\{X_i%\}_{i\in \alpha^{-1}(j)},Y_j\}_j$$ -- this is composition of multimorphisms.
%    \end{itemize}
%\end{remark}

\section{Unital $\infty-$operads}
\begin{definition}
    An $\infty-$operad $\mathcal{O}^\tensor$ is called \defn{unital} if $Mul(\emptyset,X)$ is contractible for every $X\in \mathcal{O}$.
\end{definition}

The contractible space of nullary operations will encode algebras with a unique nullary operation, for example the identity element of a monoid. It is not to be confused with the distinguished operation $id\in \mathcal{O}_1$. Requiring $\mathcal{O}$ to be unital is equivalent to requiring that it is pointed as an $\infty-$category.

%\section{Parametrized families of operads}
%We can extend the definition of operads to a functor onto $C\times FinSet_*$ that is an operad over each $c\in C$. This gives a \defn{family of $C-$operads}, and is equivalent to a generalized operad $\mathcal{O}$ with space of unary operations $\mathcal{O}_{\langle 0 \rangle}\iso C$.
%\section{The little cubes operad}
\section{Algebras over an operad}
An algebra over $O^\tensor\rightarrow FinSet_*$ valued in a symmetric monoidal category $C^\tensor\rightarrow FinSet_*$ is just a functor making $O^\tensor\rightarrow C^\tensor$ making the diagram commute and preserving inert morphisms (those that lie over inerts in $FinSet_*$.)
\end{comment}

\begin{comment}\subsection{Graph operads} The category $\mathbf{Graph}$ of simple graphs (possibly with loops and parallel morphisms) where we allow edges to contract is a perfect operator category. The terminal object is the one-vertex graph, and fibers are full subgraphs over points. $(0-1,1)$ is a point-classifier, and the right adjoint is given by adjoining a unique vertex (sent to $0$) with an edge to every other vertex. 
%Points are loops in your graph (this is already very odd). Fibers of $f:C\rightarrow D$ over a loop $\gamma$ are given by the subgraph of $C$ sent to the $\gamma$ by $f$. Let $\beta$ be the graph
%\begin{tikzcd}
%a \arrow[r, no head] \arrow["0"', no head, loop, distance=2em, in=215, out=145] & b \arrow["1"', no head, loop, distance=2em, in=35, out=325]
%\end{tikzcd}. Then the pair $(\beta,1)$ is a point classifier. The right adjoint to the functor $\mathbf{Graph}/\beta\rightarrow \mathbf{Graph}$ takes a graph $C$ to the graph $\hat{C}$ given by adding to $C$ a unique vertex $*$ with an edge to every vertex (including itself), and the unique map $\hat{C}\rightarrow \beta$ that sends all vertices of the subgraph $C$ to $a$ and $*$ to $b.$
The Leinster category is the category of finite simple graphs and partially defined maps.

\section{Cyclic operads}
Connes cyclic category is a subcategory of $Cyc$ where only degree one maps are considered - that is the corresponding map $f:S^1\rightarrow S^1$ has degree 1. It has the property that every map $\langle n \rangle \rightarrow \langle m \rangle$ factorises as a cyclic permutation of order $n+1$ followed by a map in $\Delta$. Could this be the inert-active factorisation system of a Leinster category?

Unfortunately, not. 7.2.2 of Barwick says that the isomorphisms are exactly the inert and active morphisms. All cyclic permutations are isomorphisms, but in general they are not members of $\Delta$. We don't even get a WFS, let alone an orthogonal one.

\section{Aside}
\subsection{More on loop spaces}
Consider a loop space $\Omega X$, and and choose two loops $\gamma,\delta$, viewed as maps $\gamma:I\rightarrow X$ such that $\gamma(0)=\gamma(1)=x$.  we can (non-uniquely) define their concatenation $\delta\bullet\gamma$ for instance by defining
$$\big(\delta\bullet\gamma\big)(t)=\begin{cases} 
      \gamma(2t) & 0\leq t <\frac{1}{2} \\
      \delta(2t) & \frac{1}{2} \leq t \leq 1
   \end{cases}$$
This is not unital on the nose, but is unital up to homotopy with respect to the constant path $e$: $e\bullet \gamma \simeq \gamma \simeq \gamma \bullet e$. Similarly, it is not associative on the nose, but is associative up to homotopy: $\alpha \bullet (\beta \bullet \gamma)\simeq (\alpha \bullet \beta) \bullet \gamma $. Given choices of such homotopies we can contemplate the diagram
% https://tikzcd.yichuanshen.de/#N4Igdg9gJgpgziAXAbVABwnAlgFyxMJZARgBoAGAXVJADcBDAGwFcYkQAdDptAC3q4AjZo0YwcACiHj6AAiEixkrgHN6AW3UCOw0ePkdYjHPQCUpkAF9S6TLnyEUAZgrU6TVuyndGfbbqUDQRlZUwU9ZQ41TX9FfS4jEwtrW2w8AiIAFlJiNwYWNkQQCW8efnDA6RNQivFVDS0wnTicA0T6KxsQDDSHInIcvI9Czh8-WslvYOqJ+pimgPjDGGMzTtT7DJQAJlJtoYKvLjLYiNkpkNmohrMFlraVkys3GCgVeCJQADMAJwh1JADEA4CBIYgpEC-f5gmggpDbCFQgGIIFwxBORF-ZEuYGgxCZTHQ-GwvEIyiWIA
\[% https://tikzcd.yichuanshen.de/#N4Igdg9gJgpgziAXAbVABwnAlgFyxMJZARgBoAGAXVJADcBDAGwFcYkQAdDptAC3q4AjZo0YwcACiHj6AAiEixkrgHN6AW3UCOw0ePkdYjHPQCUpkAF9S6TLnyEU5UsWp0mrdlx78FeyVI6Mga6SqoaWqZ+SgZGJhbWtth4BEQATKRpbgwsbIgggT7aofqBgsHR4uGaZlE6ivpccfRWNiAYyQ7pFNkeeQXejHzFDTghwXUlyhxqNZVjTTDGZq1J9qkoAMwuvbnsEoVDvvX+4yayk6PVkfOxSyZWbjBQKvBEoABmAE4Q6kjOIBwECQxESIG+vxBNCBSDSYIhf0QAJhiE28J+iO2gOBiAALOjIYgMtikPjKJYgA
\begin{tikzcd}
                                                               & \alpha\bullet(\beta \bullet(\gamma\bullet \delta)) \arrow[ld] \arrow[r] & (\alpha\bullet \beta )\bullet(\gamma\bullet \delta) \arrow[rd] &                                                     \\
\alpha\bullet((\beta \bullet\gamma)\bullet \delta) \arrow[rrd] &                                                                         &                                                                & ((\alpha\bullet \beta )\bullet\gamma)\bullet \delta \\
                                                               &                                                                         & (\alpha\bullet (\beta \bullet\gamma))\bullet \delta \arrow[ru] &                                                    
\end{tikzcd}\] and ask whether it commutes. There is no guarantee it will commute on the nose, but it will commute up to a homotopy of homotopies, and so on ad infinitum. Up to formulation, it seems loop spaces are topological monoids "up-to-homotopy", also called $E^1-$spaces. For double loop spaces $\Omega^2 X$ we have two operations: given two loops $\gamma,\delta:I\rightarrow \Omega X$ we can either concatenate "externally" as before
$$\big(\delta\bullet_0 \gamma\big)(t)=\begin{cases} 
      \gamma(2t) & 0\leq t <\frac{1}{2} \\
      \delta(2t) & \frac{1}{2} \leq t \leq 1
   \end{cases}$$
or we can concatenate "internally":
$$\big(\delta\bullet_1 \gamma\big)(t)=\delta(t)\bullet \gamma(t).$$
These two operations have the same unit (up to homotopy) and distribute over each other\footnote{This is maybe not clear immediately.}. By an Eckmann–Hilton argument this should correspond to a single commutative product on $\Omega^2X$ - but what about the homotopies? As we will see later, the product commutes up to homotopy, but these homotopies do not commute up to homotopies. This is called an $E_2$-space. Continuing this process, May showed that n-fold loop spaces are $E_n$-spaces, including the case $n=\infty$ where the product maximally commutes. Furthermore, this is an if and only if statement: a space $Y$ is the homotopy type of an n-fold loop space $\Omega^nX$ if and only if $Y$ is an $E_n$-space. This is May's recognition principle. In order to state it fully, we need a more manageable definition of an $E_n$-space. Operads give a way to do this, furthermore for any algebraic structure in any symmetric monoidal category.

Operads correspond to monads functorially by mapping $\mathcal{O}$ to  $X\mapsto \int^j \mathcal{O}(j)\times X^j$. We will soon show that we have operads $C_n$ with weak homotopy equivalences $C_nX\rightarrow \Omega^n\Sigma^n X$ when $X$ is connected - this will be a statement about monads.

We will work in a convenient monoidal closed category of topological spaces.

\begin{definition}
    An $A_\infty$ operad $C$ is a $\Sigma$-free operad over $Ass$ (there is a map that's an equivalence on $\pi_0 C\rightarrow Ass$) such that $C\rightarrow M$ is a local $\Sigma$-equivalence (each $C(j)$ is $\Sigma$-equivarently homotopy equivalent to $Ass(j)$).

    An $E_\infty$ operad $D$ is a $\Sigma-$free operad over $Com$ such that $D\rightarrow Com$ is a local equivalence.
\end{definition}

In other words, an $E_\infty$ is an operad $C$ such that each $C(j)$ is $\Sigma_j$-free and contractible.

\begin{definition}
    Let $C_n$ be the little cubes operad. We have $C_n(j)$ is the space of $j$ pairwise disjoint embeddings of open $n-$cubes into $\partial I^n$. There is an obvious composition, unit and $\Sigma$-action. 
\end{definition}

Note that $C_1$ is $A_\infty$ and $C_\infty$ is $E_\infty$.

The recognition principle is this: Every $\Omega^nX$ is a $C_n$ space and every connected $C_n$ space $C_nX$ is weakly equivalent to $\Omega^n \Sigma^n X$.

\begin{example}
    Consider a convenient\footnote{compactly generated Hausdorff} category of topological spaces $(Top,\times,*)$. The \defn{little $n-$cubes operad} $C_n$ has as $k^{\text{th}}$ space of operations $C_n(k)$ the space of rectilinear disjoint embeddings of $k$ open $n-$cubes inside an $n-$cube.\footnote{This is topologized with the subspace topology $Rect(Cube_n\times k,Cube_n)\subset (\mathbb{R}^{2n})^k$.} A rectilinear embedding is given by a product of affine functions. The identity embedding is the identity. Composition is given by pasting (show picture, remember numbering). The symmetric action permutes cubes.   

    An element of $C_n(k)$ parametrizes a way to compose $k$ $n-$loops in a loop space $\Omega^n X$. Associativity and unitality is now globally kept track of. The continuous paths correspond to homotopies of compositions, indeed a continuous map

    $$I\rightarrow C_n(k)\rightarrow Top((\Omega^nX)^k,\Omega^nX)$$
    corresponds to a continuous map $$I\times (\Omega^nX)^k\rightarrow \Omega^nX$$ i.e. a homotopy, since $\times$ is closed in the convenient category. Similarly, a continuous map $I^2\rightarrow C_n(k)$ (rel endpoints) corresponds to a homotopy of homotopies.

    The topology of $C_n(k)$ makes it clear why $E_n$-spaces are algebras. For $k=1$ it is clear that we cannot commute up to homotopy. At $k=2$ it is now clear that we can commute, but consider the homotopically non-trivial path from a fixed $f\in C_2(2)$ to itself. This is non-trivial for $\pi_1S^1$ reasons, or if you like, braiding reasons. However, any two braidings in a 3-cube can be continuously deformed to each other. But notice the homotopically non-trivial homotopy between the identity homotopy on the identity path on $f\in C_3(2)$. This is for $S^2$ reasons, etc.
    
    May's recognition principle is that $E_n$-spaces are exactly homotopy types of n-fold loop spaces.
\end{example}
\end{comment}


%%%


%\begin{example}
%Let $\mathbf{Cyc}$ be the category of finite cyclically ordered sets and monotone functions. A cyclic order on a set $X$ is a ternary relation such that
%\begin{itemize}
%    \item $[a,b,c]\implies [c,a,b]$
%    \item $[a,b,c]\implies ![c,b,a]$
%    \item $\big([a,b,c]$ and $[a,c,d]\big)\implies [b,c,d]$
%    \item if $a,b,c$ are distinct then either $[a,b,c]$ or $[a,c,b]$.
%\end{itemize}
%This is equivalent to the category of homotopy classes of maps $F\rightarrow S^1$ from finite sets such that $f(x)\neq f(y)$ if $x\neq y$, where maps are maps of finite sets such that the obvious triangle commutes. (Notice this is similar to the world-object of $S^1$ but only spanned by separated decoupling structures ).
%A function $f:X\rightarrow Y$ is \defn{monotone} if and only if $[f(x),f(y),f(z)]\implies [x,y,z]$. Then $\mathbf{Cyc}$ is an operator category with the obvious terminal object. The fiber of a map $f:I\rightarrow J$ at $j\in J$ is the subset $f^{-1}(j)$ given the \textit{subset order}. There is a universal point $(2=\{0,1\},0)$ where $\{0,1\}$ is given the unique (empty) order. This is the only universal point up to isomorphism. However, $\mathbf{Cyc}$ is not perfect. Indeed, a monotone map $I\rightarrow 2$ is just a function, which we can think of a 2-coloring of $I.$ For every $C'\subset C\in \mathbf{Cyc}$, let $C_{C'}$ denote the map $C\rightarrow 2$ whose fiber over $0$ is $C_1$. Now consider $H\in\mathbf{Cyc}$ and note there is no functor $\hat{(-)}:\mathbf{Cyc}\rightarrow \mathbf{Cyc}/2$ with natural bijections $$\text{Hom}(C_{\emptyset},\hat{H}_{H'})\iso \text{Hom}(\emptyset,H)\iso *$$ and
%$$\text{Hom}(H_H,\hat{H}_{H'})\iso \text{Hom}(H,H').$$
%Indeed, together these conditions require that $H'\iso H$ while $\hat{H}\setminus H' \iso *$. Clearly there is not a universal way to add a single point to a cyclically ordered set.
%We claim $\mathbf{C}$ is perfect. The right adjoint is given by adding points inbetween every pair - i.e. $E(I)=I\sqcup I$ with $[i,j,k]\implies [(i,0),(i,1),(j,0)]$ and $[(j,0),(j,1),(k,0)]$. The special map $E(I)\rightarrow 2$ is given by $(i,j)\mapsto j.$ The counit of the adjunction is the identity. For the unit, note a map $I\rightarrow 2$ is the same as a coloring of $I$ in red and blue. Then the unit $I\rightarrow E(I)$ is the unique map which restricts to the inclusion in the first component on $f^{-1}(0).$
%PROBLEM: This is not an adjunction:( Take $C$ with the entirely blue coloring. Then there is a unique map $\empty\rightarrow H$, but not a a unique map from $C$ to $H$.
%An operad over $C$ is...
%The terminal operad in $C$ is the little $1-$cubes on $S^1$ operad. The terminal operad in $\prod_n S^1$ is the little $n-cubes$ on $\prod_n S^1$ operad. OR IS IT? What would composition look like?
%\end{example}